\chapter{Mapa das posições escatológicas (convergências, diferenças e por que isso importa menos do que você pensa)}
\noindent\textit{Vinheta.} Dois líderes de célula discutem sobre o ``milênio'' no WhatsApp. Em 10 minutos, já não se escutam; trocam rótulos. No domingo, um visitante novo pergunta: ``\emph{Jesus volta mesmo?}'' Os dois concordam na hora: ``\textbf{Sim}.'' Às vezes esquecemos o \textbf{essencial} porque brigamos pelos \textbf{detalhes}.

\section{O que nos une (núcleo ortodoxo)}
\begin{checklist}
\begin{itemize}[leftmargin=*,noitemsep]
  \item \textbf{Jesus voltará} de forma \textbf{pessoal, visível e gloriosa} (At~1:11; Mt~24:30).
  \item \textbf{Ressurreição} e \textbf{juízo} final (Jo~5:28--29; Ap~20:11--15).
  \item \textbf{Vitória definitiva de Cristo} sobre o mal e \textbf{Nova Criação} (Ap~21--22).
  \item \textbf{Chamado à santidade e missão} enquanto esperamos (2~Pe~3:11--12).
\end{itemize}
\end{checklist}

\section{Quatro lentes de leitura (escopos históricos)}
\begin{itemize}[leftmargin=*,noitemsep]
  \item \textbf{Preterista}: muitos cumprimentos no contexto do séc.~I (Roma, igrejas da Ásia).
  \item \textbf{Historicista}: o livro percorre a história da igreja em \emph{eras}.
  \item \textbf{Futurista}: ênfase no cumprimento literal próximo ao fim.
  \item \textbf{Idealista}: princípios recorrentes válidos para todas as gerações.
\end{itemize}
\noindent\textit{Integração sugerida neste livro}: usar o \textbf{melhor de cada} (cap.~1) sem absolutizar uma lente.

\section{Três posições sobre o milênio (Ap 20)}
\subsection*{1) Amilenismo (milênio como símbolo do presente século)}
\textbf{Resumo}: ``Milênio'' = período entre a primeira e a segunda vinda; Cristo reina agora; at amarração de Satanás = limitação do seu engano às nações (Mc~3:27).\\
\textbf{Força}: enfatiza a \textbf{centralidade de Cristo agora} e a leitura simbólica coerente.\\
\textbf{Risco}: minimizar expectativas de futuros juízos/conflitos históricos concretos se lido de modo \emph{apenas} idealista.

\subsection*{2) Premilenismo (Cristo volta antes do milênio)}
\textbf{Resumo}: Jesus retorna, \textbf{vence} e estabelece um reino milenar na terra; depois vem o juízo final.\\
\textbf{Variante histórica}: foco em continuidade igreja-história; tribulação não necessariamente \emph{pré}-arrebatamento.\\
\textbf{Variante dispensacional}: distingue fortemente Israel/igreja e inclui \emph{arrebatamento} anterior a uma grande tribulação.\footnote{Este livro evita esquemas fechados; onde o texto for explícito, afirmamos; onde for silencioso, mantemos humildade.}\\
\textbf{Força}: toma seriamente as \textbf{promessas concretas} de restauração e a \textbf{esperança futura visível}.\\
\textbf{Risco}: leituras jornalísticas forçadas; cronogramas rígidos não exigidos pelo texto.

\subsection*{3) Pós-milenismo (o evangelho avança na história)}
\textbf{Resumo}: o reino de Deus progride historicamente; ao final, Cristo volta; ``milênio'' pode ser longo período de expansão do evangelho.\\
\textbf{Força}: destaca \textbf{otimismo do evangelho} e cultura redimida.\\
\textbf{Risco}: confundir progresso cultural com \textbf{consumação}; subestimar a persistência do mal.

\section{Questões em debate honesto}
\begin{itemize}[leftmargin=*,noitemsep]
  \item \textbf{Israel e igreja}: continuidade/aliança ou distinções específicas no plano histórico?
  \item \textbf{Tribulação}: momento, duração e relação com a igreja.
  \item \textbf{Sequência de Ap 19--21}: quanto é \textbf{cronologia} e quanto é \textbf{cena} (como nos evangelhos)?
  \item \textbf{Leitura de números}: literais ou simbólicos (``mil'', ``1260 dias'', ``sete'')?
\end{itemize}

\section{NIC aplicado ao mapa}
\begin{nicbox}
\textbf{Narrativa}: toda posição responsável deve contar a \textbf{boa nova} do Cordeiro que vence.\\
\textbf{Identidade}: qualquer esquema que \textbf{apague} a santidade e a missão presentes está torto.\\
\textbf{Controle}: cuidado com sistemas que viram \textbf{controle tribal} (``nós certos x eles errados'').
\end{nicbox}

\section{Como conviver e ensinar (pacto de caridade)}
\begin{checklist}[title={Regras simples para igrejas saudáveis}]
\begin{itemize}[leftmargin=*,noitemsep]
  \item \textbf{Cristo no centro}: pregue textos, não \textbf{timelines}.
  \item \textbf{Essenciais vs. adiáforas}: firme nos \textbf{essenciais}; caridoso nos secundários.
  \item \textbf{Marcar especulação}: frases com ``\emph{parece}'' e ``\emph{pode ser}'' quando extrapolar.
  \item \textbf{Pluralidade segura}: permita classes/leituras diferentes sob guarda de credo confessional comum.
  \item \textbf{Prática antes de polêmica}: aplique esperança, santidade e missão toda semana.
\end{itemize}
\end{checklist}

\section{Ferramenta — Comparador em 10 perguntas (grupo de estudo)}
Use com qualquer posição; responda em 1--2 frases cada.
\begin{checklist}
\begin{itemize}[leftmargin=*,noitemsep]
  \item Quem reina \textbf{agora} e \textbf{como}? (Ap~5; Mt~28:18)
  \item O que significa a \textbf{amarra} de Satanás? (Ap~20:2)
  \item O ``\textbf{milênio}'' é literal ou simbólico? Por quê?
  \item Como sua posição lê \textbf{Ap 19--21}: sequência cronológica ou cenas?
  \item Onde entra \textbf{Israel} no plano de Deus após Cristo? (Rm~9--11)
  \item O que você espera da \textbf{igreja} no mundo até a volta de Cristo?
  \item Como sua leitura protege contra \textbf{pânico} e contra \textbf{comodismo}?
  \item Quais \textbf{textos difíceis} para sua posição e como você os trata?
  \item O que essa posição \textbf{edifica} na prática na sua comunidade?
  \item O que você está disposto a \textbf{revisar} diante do texto?
\end{itemize}
\end{checklist}

\section{Leitura paralela (4 semanas, igreja ou grupos)}
\begin{checklist}[title={Plano simples}]
\begin{itemize}[leftmargin=*,noitemsep]
  \item \textbf{Semana 1}: Daniel 7; Marcos 13; Apocalipse 1.
  \item \textbf{Semana 2}: 2~Tessalonicenses 2; Apocalipse 12--13.
  \item \textbf{Semana 3}: Romanos 8; 1~Tessalonicenses 4--5; Apocalipse 19.
  \item \textbf{Semana 4}: Apocalipse 20--22; 2~Pedro 3.
\end{itemize}
\end{checklist}

\section{Perguntas para grupos pequenos}
\begin{itemize}[leftmargin=*,noitemsep]
  \item Em quais pontos você percebe \textbf{convergência} com irmãos de outra posição?
  \item Qual \textbf{risco} a sua tendência pessoal traz (pânico, presunção, comodismo)?
  \item Como sua comunidade pode \textbf{ensinar com caridade} sem abrir mão da convicção bíblica?
\end{itemize}

\bigskip
\noindent\textit{Próximo capítulo: Perguntas difíceis — doze objeções comuns e respostas bíblicas e práticas.}
